\documentclass[../cours_tle.tex]{subfiles}
\begin{document}
\setcounter{chapter}{15}
\chapter{Logarithme, puissances et modèles multiplicatifs}
\label{ch:16}
\begin{objectif}
Spécialité. Construire le logarithme et justifier ses règles. Prérequis : exponentielle, continuité et dérivation.
\end{objectif}
\section{Inverser l'exponentielle}
L'exponentielle est continue, strictement croissante et a pour limites $0$ et $+\infty$ aux deux infinis. Le théorème des valeurs intermédiaires montre qu'elle est une bijection de $\R$ sur $]0,+\infty[$.
\begin{definition}
Pour $x>0$, $\ln x$ est l'unique réel $y$ tel que $e^y=x$.
Ainsi $\ln(e^y)=y$, $e^{\ln x}=x$, $\ln1=0$ et $\ln e=1$.
\end{definition}
Le logarithme est croissant et continu comme réciproque de cette bijection (chapitre 13). Les courbes de $\ln$ et de $\exp$ sont symétriques par rapport à $y=x$ car leurs coordonnées s'échangent.
\section{Propriétés algébriques, sans règles inventées}
Pour $a,b>0$, $e^{\ln a+\ln b}=ab$. Par unicité de l'antécédent,
\[
\ln(ab)=\ln a+\ln b,\quad \ln(a/b)=\ln a-\ln b,\quad \ln(a^n)=n\ln a\ (n\in\Z).
\]
La formule du quotient suit de $\ln(1/b)=-\ln b$ ; celle des puissances suit par récurrence puis par passage à l'inverse.
Il n'existe pas de règle $\ln(a+b)=\ln a+\ln b$ : avec $a=b=1$, elle donnerait $\ln2=0$.
\section{Dérivée, avec existence démontrée}
\begin{theoreme}
Pour $x>0$, $(\ln x)'=1/x$. Pour $u>0$ dérivable, $(\ln u)'=u'/u$.
\end{theoreme}
\begin{demonstration}
Fixer $a>0$ et poser $k=\ln(a+h)-\ln a$. La continuité donne $k\to0$, et $h=a(e^k-1)$. Pour $h\ne0$, $k\ne0$, donc
$\dfrac{\ln(a+h)-\ln a}{h}=\dfrac1a\dfrac{k}{e^k-1}\to1/a$, puisque $\exp'(0)=1$.
La règle de composition donne la seconde formule.
\end{demonstration}
La dérivée est positive ; la dérivée seconde vaut $-1/x^2$, donc $\ln$ est concave.
Sur un intervalle où $u$ ne s'annule pas, $(\ln|u|)'=u'/u$, en séparant les signes de $u$.
\section{Limites et croissances comparées}
Pour tout $A$, $x>e^A$ entraîne $\ln x>A$, donc $\ln x\to+\infty$.
Pour $x\to0^+$, utiliser $\ln x=-\ln(1/x)$ : la limite est $-\infty$.
Pour $m>0$, poser $t=\ln x$ : $(\ln x)/x^m=t/e^{mt}\to0$ quand $x\to+\infty$.
À $0^+$, poser $x=e^{-t}$ : $x^m\ln x=-te^{-mt}\to0$. En particulier $x\ln x\to0$, malgré $\ln x\to-\infty$.
\section{Équations, seuils et puissances réelles}
$\ln(x-1)+\ln(x+1)=\ln3$ impose $x>1$. On obtient $x^2-1=3$, donc seulement $x=2$.
Pour $0<q<1$, $q^n<\varepsilon$ équivaut à $n>\ln\varepsilon/\ln q$ : le sens se renverse car $\ln q<0$.

En complément, pour $x>0$ et $\alpha\in\R$, définir $x^\alpha=e^{\alpha\ln x}$ donne
$(x^\alpha)'=\alpha x^{\alpha-1}$ et les règles usuelles de puissances par celles de l'exponentielle.
La définition exclut les bases négatives pour un exposant réel arbitraire.
\section{Briggs et une limite remarquable}
Prendre plusieurs racines carrées rapproche $x>0$ de $1$. Si $r_k=x^{1/2^k}$, alors
$2^k\ln r_k=\ln x$ et $\ln(1+h)/h\to1$.
Ainsi $2^k(r_k-1)\to\ln x$ : c'est une forme de la méthode de Briggs. Trop de racines provoquent une perte de précision en flottants.
Enfin $n\ln(1+x/n)\to x$ pour tout réel fixé $x$ ; par continuité de l'exponentielle, $(1+x/n)^n\to e^x$, les bases étant positives pour les grands $n$.
\begin{finchapitre}
\section{Utilité et suite du parcours}
Le logarithme transforme des produits en sommes : croissance composée, ordres de grandeur, vraisemblance statistique, entropie et analyse des algorithmes. Sa dérivée explique son lien avec les intégrales de 1/x.
\begin{retenir}
Toujours imposer un argument positif. Lors d'une division par un logarithme négatif, inverser le sens de l'inégalité.
\end{retenir}
\end{finchapitre}
\end{document}
